% compendium/execution-trading-costs/dark-trading-fragmentation.tex
% Entries migrated on 2026-10-16 from the legacy hidden-liquidity-and-dark-trading and cross-venue-routing-and-
% fragmentation files (slice 34), owned by market-structure-execution-quality, a map candidate of both files, whose
% outcome names fragmentation and dark trading. Other blocks of the same files were merged into existing entries
% (order exposure, block trades, smart order routing, fee-adjusted venue choice). Dispositions in
% notes/migration-slices.md; the files verbatim in notes/artefacts/market-structure-execution-quality/.

\section{Dark trading and fragmentation}

\begin{qfdefinition}[Dark pools and crossing networks: derivative pricing, the kinds of dark venue, integrated and midpoint books, why traders use them, what the evidence says, and their risks]\label{def:dark-pools-crossing}
\qffield{Statement} A \emph{dark pool} is a venue that does not publicly display its available liquidity; its main interest is to protect institutions executing large orders from opportunistic short-horizon traders \citep[chap.~1, sec.~1.4, p.~102]{alpha-laruelle}. Dark trading existed before as internal crossing, upstairs trading and OTC trading (p.~103). Lehalle and Laruelle list, for European markets after MiFID (pp.~103--104): \emph{crossing networks}, which cross buy and sell orders at a price, typically the mid-price, derived from the primary market; \emph{internalisation pools}, in which a broker crosses its own flow; \emph{electronic market makers}, which trade against the operator's flow and accept or reject incoming orders at once; \emph{exchange-based pools}, hidden liquidity on lit exchanges executed only against sufficiently large contra orders; \emph{dark aggregators}, which split an order across pools and adjust as it fills; and \emph{size-discovery} (large-block) pools created in response to MiFID~II's caps. Crossing networks are the only order-driven markets that are not auctions: all trades take place at prices determined elsewhere, under a \emph{derivative pricing} rule, so they discover only whether traders will trade at those prices, not the prices themselves \citep[chap.~6, sec.~6.5, p.~132]{alpha-harris}.

\qffield{Notation} Crossing rules: Harris's 2002 U.S.\ networks were calls; Instinet's after-hours cross and the NYSE's After-hours Session~I crossed at the 4 p.m.\ close, and POSIT crossed eight times a day at the primary midpoint at a random time within seven minutes after each call, allocating excess supply or demand pro rata; some networks operate continuously (pp.~132--134). The main European MTFs' dark books, under MiFID's reference-price waiver, crossed only at the midpoint of a reliable reference market's best bid and offer, with a two-book design: an \emph{integrated book} for displayed orders and hidden orders large in scale (LIS), and a \emph{dark midpoint book} for smaller hidden orders, which match only at the primary midpoint; they accepted mid-pegged dark orders with limit price, minimum acceptable quantity and immediate-or-cancel options (Lehalle and Laruelle, p.~104).

\qffield{Conditions} \emph{Derivative prices} must be neither stale nor manipulated \citep[chap.~6, sec.~6.5.2, pp.~134--136]{alpha-harris}. Informed traders trade eagerly at a stale price and refuse otherwise, an adverse-selection problem: the after-hours crosses received more buy orders when prices had risen after the close, and did not cross stocks with significant after-hours news (p.~135). Anyone who will trade at a price fixed elsewhere later may try to move it: Harris's example saves \$15{,}000 on 500{,}000 shares by depressing the reference price 3 cents with a 1{,}000-share sale (p.~135). \emph{Price ownership.} Primary markets argue that crossing networks skim the easiest orders and use their prices without paying for them; crossing customers reply that they should not pay for price discovery in which they do not take part, and that their own orders in the primary market produce those prices (sec.~6.5.1, p.~134). \emph{Price discovery.} Lehalle and Laruelle state that dark pools cannot create prices and influence price formation indirectly by attracting activity; the literature they cite disagrees, some empirical studies finding reduced price discovery and others, including Zhu's model, improved informativeness of exchanges (pp.~105--107); Zhu (\texttt{wide-zhu-dark-pools}) and Buti et al.\ are not read here.

\qffield{Practice} \emph{Why use them.} Secrecy, to reduce information leakage and impact, and price improvement, trading at the midpoint without paying the whole spread (Lehalle and Laruelle, p.~108). \emph{Evidence, as reported.} From Buti, Rindi and Werner's U.S.\ study, as Lehalle and Laruelle summarise it: dark activity is concentrated in liquid stocks; for a given stock it is higher on days with high volume, higher depth and lower intraday volatility, and lower on days with large order imbalances relative to volume; it is associated with better spreads, depth and volatility and with short-term overreaction (p.~107). \emph{Checked against the working paper} (2026-10-16, the Dice Center version of 24 June 2010, from an Internet Archive copy of Ohio State's working-paper page; the 2022 \emph{Financial Management} version was not read). The data are daily single-counted share volume that eleven dark pools reported voluntarily to SIFMA for calendar 2009, about a third of the roughly 32 active U.S.\ pools, anonymous by pool, with no public data to check them against and no way to know whether the eleven are a stable fraction across stocks and days \citep[pp.~5--6]{brw-2010}. Dark activity is concentrated in large, high-volume, high-price stocks with low spreads, high depth and low short-term volatility, and on a stock's days with high volume and depth and low volatility, and is lower on days with large order imbalances relative to volume and large absolute returns (abstract, p.~2). On market quality the authors stress that dark activity and market quality are ``jointly determined'' and claim at best an association: cross-sectionally, more dark activity goes with lower spreads, price impact and short-term volatility; in a matched-pair design (100 NASDAQ and 100 NYSE stocks, matched on size and price, dark activity measured January--June 2009 and market quality July--December), high-dark-activity stocks had quoted spreads 3.51~bp and effective half-spreads 0.53~bp lower (sec.~6, pp.~13--16). This draft's price-efficiency section is marked ``TBA'' (p.~6), so the ``short-term overreaction'' finding Lehalle and Laruelle report could not be checked in it. \emph{Theory} (Zhu, read in the published version on the author's page). In a model where informed and liquidity traders both choose between an exchange and a midpoint dark pool without guaranteed execution, informed orders cluster on the heavy side and fill less often in the dark, so the dark pool attracts relatively more uninformed flow; ``under natural conditions'' this raises the exchange's signal-to-noise ratio and lowers pricing error while widening the exchange spread, so better price discovery comes with lower exchange liquidity; the result needs short-lived information (it weakens as the information horizon lengthens), abstracts from pinging, routing and indications of interest, and makes no welfare claim \citep[pp.~747--750, Propositions~3--4, pp.~763 and~770]{wide-zhu-dark-pools}. In his numerical example adding a dark pool cuts the market maker's percentage pricing error from 7.1 to 2.6 percent at high volatility (p.~770). Zhu's result is a model's, and Buti, Rindi and Werner's an association in one year's partial data; neither measures price discovery in current markets. Lehalle and Laruelle's own one-minute data for CAC~40 and FTSE~100 stocks on Chi-X, BATS and Turquoise dark books in the first three weeks of March 2010 show dark \emph{turnover} concentrated in the most volatile deciles (nearly a quarter of it in the top tenth), but dark \emph{market share} not rising with volatility, and rising strongly with the spread (about three times as high in the widest-spread decile as in the narrowest), most of all for less liquid SBF~120 and FTSE~250 stocks (pp.~108--114). \emph{Fill rates.} Harris reports that U.S.\ crossing networks filled less than 10 percent of the order volume submitted, and that unfilled orders were usually sent elsewhere (p.~134); a dark fill rate measured only from executions therefore misses the orders that found nothing. \emph{Bilateral and multilateral.} On a multilateral platform anyone may supply or take liquidity; bilateral trading keeps an asymmetry that protects liquidity providers from toxic flow and is easier to sustain in illiquid instruments (Lehalle and Laruelle, p.~114). \emph{Broker crossing networks} (slice~35). Lehalle and Laruelle describe executing brokers, squeezed by the fixed costs of fragmentation and falling assets after the crisis, opening pools marketed as ``100\% first-hand liquidity''; these BCNs sometimes also hosted proprietary flow of the owning bank, MTFs and some exchanges launched dark pools in response, and some brokers converted BCNs into MTF dark pools to reassure clients \citep[chap.~1, sec.~1.1.4, pp.~55--56]{alpha-laruelle}; large investors favoured BCNs ``as far as such pools are `HFT-free guaranteed' '' (Introduction, p.~18). The legacy taxonomy (dark pools hide pre-trade interest; internalisers match against a broker-dealer's inventory or flow; crossing networks match natural interest at reference prices) restates this entry; the current EU caps are in \hyperref[rule:eu-trading-venue-transparency]{EU trading venues and transparency}.

\qffield{Pitfalls} \emph{Gaming and toxicity.} Some high-frequency traders detect large dark orders by sending small orders (``fishing'') or orders on the opposite side, then move the primary price; protections include large minimum quantities, restricted access, tracking participants who behave like gamers and statistical detection of consistently profitable participants; indications of interest can leak information too (Lehalle and Laruelle, p.~115). \emph{Internal crossing.} Brokers who cross internally may expose client orders only to their own clients, which makes the search for counterparties harder, and may favour some clients over others; brokers reply that they work harder for two commissions, and exposure to the market protects clients (Harris, chap.~25, sec.~25.4, pp.~520--521; \hyperref[def:broker-agency-conflicts]{broker conflicts}). The legacy finding that dark activity is concentrated in ``volatile periods'' holds for dark \emph{turnover} in Lehalle and Laruelle's deciles but not for dark \emph{market share}, and conflicts with the lower-volatility days in Buti et al.'s results as reported; it is kept only in that qualified form. Venue lists, market shares (about 5 percent of non-OTC European trades and 10--12 percent of U.S.\ share volume, p.~105), the LIS thresholds and the MiFID~II caps are dated and not current; the 2024 EU review is a backlog item of this pathway.

\qffield{See also} \hyperref[mech:order-matching-pricing]{matching and trade pricing} (derivative pricing and pro-rata allocation); \hyperref[def:order-exposure-decisions]{order exposure decisions} (hidden orders and probing); \hyperref[def:block-trade]{block trades}; \hyperref[def:smart-order-routing]{smart order routing}; \hyperref[def:best-execution]{best execution}; \hyperref[def:broker-agency-conflicts]{broker conflicts}.

\qffield{Further reading} Zhu (published 2014) and Buti, Rindi and Werner (2010 working paper) were read for this entry (Practice); the latter's published version, its later price-efficiency results and the Rosenblatt dark-pool reports are not read. \emph{Empirical evidence on price discovery} (batch of 2026-10-17). Comerton-Forde and Putniņš use ASX data from February 2008 to October 2011 and find that ``dark trades are less informed than lit trades'' and that high levels of dark trading raise adverse selection on the lit exchange. They find ``low levels of non-block dark trading are benign or even beneficial for informational efficiency, but high levels are harmful'', with deterioration from about 10 percent of a stock's dollar volume, while ``no evidence that block trades in the dark impede price discovery'' is found; block trades away from the lit market up to about 40 percent of dollar volume can help \citep[author's accepted manuscript of 13 February 2015, abstract and PDF pp.~4--5 and~10]{comerton-forde-putnins-2015}. The thresholds are for one market and period, before the EU's double volume caps and their later revision. The legacy file is kept verbatim in \nolinkurl{hidden-liquidity-dark-trading.tex} in the folder \nolinkurl{notes/artefacts/market-structure-execution-quality/}.

\qffield{Index} dark pool; crossing network; derivative pricing rule; reference price; midpoint cross; POSIT; internalisation pool; electronic market maker; exchange-based pool; dark aggregator; size-discovery pool; large in scale; reference-price waiver; integrated book; midpoint book; stale price; price manipulation; price ownership; dark market share; fishing; gaming; toxic flow; bilateral trading; multilateral trading; internal crossing; broker crossing networks; first-hand liquidity

\qffield{Status} Verification: partly checked 2026-10-16 (against the text layers of local copies of \texttt{alpha-laruelle}, 2018 World Scientific printing, chap.~1, sec.~1.4, pp.~102--115, and \texttt{alpha-harris}, 2002 Oxford printing, chap.~6, pp.~132--136, and chap.~25, pp.~520--521; slice~35: pp.~18 and 55--56; source follow-up of 2026-10-16: \texttt{wide-zhu-dark-pools}, pp.~747--770, and Buti, Rindi and Werner's June 2010 working paper, pp.~2--16; batch of 2026-10-17: Comerton-Forde and Putniņš, accepted manuscript from the UTS repository, abstract and PDF pp.~4--5 and~10, text layer). Stability: draft. Currency: U.S.\ crossing networks as of 2002; U.S.\ dark pools in 2009; European venues and data as of 2010--2018; no current rule stated.
\end{qfdefinition}

\begin{qfmethod}[Venue market share: value and trade-count shares, auctions and the ex-fixing share, aggregation across stocks, estimation error, and why shares depend on the time of day]\label{meth:venue-market-share}
\qffield{Statement} For stock $k$ traded on $N$ venues between $t_1$ and $t_2$, with trade $\ell$ at price $P^k_\ell$, size $V^k_\ell$, time $\tau^k_\ell$ and venue $\delta^k_\ell$, Lehalle and Laruelle define the share of venue $n$ in traded value and in trades \citep[chap.~1, sec.~1.1.1, eqs.~(1.1)--(1.2), pp.~34--35, page image of p.~35]{alpha-laruelle}:
\[ \mathfrak M^k(n)=\frac{\sum_{t_1\le\tau^k_\ell\le t_2}P^k_\ell V^k_\ell\mathbf 1_n(\delta^k_\ell)}{\sum_{t_1\le\tau^k_\ell\le t_2}P^k_\ell V^k_\ell},\qquad \mathfrak M^k_T(n)=\frac{\sum_{1\le\ell\le T}\mathbf 1_n(\delta^k_\ell)}{T}, \]
the value share being close to the share in shares while the price varies little. Across stocks the share can be weighted by turnover, by index weight (free float or capitalisation, not turnover) or equally, and the answers differ (pp.~35--37). The right metric depends on the question (p.~37): the all-included share for a venue's position and revenue; for the probability that a \emph{trade} occurs on venue $n$, the trade-count share, with auctions and continuous trading measured separately; for the probability that a unit of \emph{currency} trades there, the value share, again separately for auctions. Treating market share as a proportion, its estimator for asset $k$ and venue $i$ is the turnover on $i$ over total turnover, $\hat p^i_k$, and an index-level share is $\hat p^i=\sum_kw_k\hat p^i_k$ (app.~A, secs.~A.8.1--A.8.2, pp.~286--287).

\qffield{Notation} $m(n)$ is venue $n$'s share of auction volume and $p$ the auctions' share of all volume. Interpretation, derived: venue $n$'s continuous-session share is then
\[ \mathfrak M(n;\text{ex-fixing})=\frac{\mathfrak M(n)-m(n)\,p}{1-p}, \]
since its continuous turnover is its total less its auction turnover; Lehalle and Laruelle print $(1-m(n)p)\mathfrak M(n)/(1-p)$ (p.~35, page image); the two agree exactly when $m(n)\,p\,(1-\mathfrak M(n))=0$: when the venue has no auction volume, when there are no auctions, or when the venue has all the trading (a monopolist, share 1 under both); the domain is $0\le p<1$, $0\le m(n)\le1$ and $m(n)p\le\mathfrak M(n)\le m(n)p+1-p$, and with $p=1$ the continuous share is undefined (the slice-34 wording ``agrees only when $m(n)=0$'' was too strong). Checked (\nolinkurl{slice34-checks.py}, with assertions): a primary market with 60 percent of all trading and all the auctions, with auctions 20 percent of volume, has 50 percent of continuous trading, where the printed formula gives 60; the monopolist and no-auction cases agree.

\qffield{Conditions} \emph{Estimation error.} Lehalle and Laruelle treat trades as independent Bernoulli trials, so $\operatorname{Var}\hat p=p(1-p)/n$ with a normal confidence interval $\hat p\pm z_{1-\alpha/2}\sqrt{\hat p(1-\hat p)/n}$, adequate when $np>5$ and $n(1-p)>5$ (sec.~A.8, pp.~284--286), and assuming independence across stocks they derive $\operatorname{Var}(\hat p^i_k-\hat p^i)=(1-2w_k)p^i_k(1-p^i_k)/n^i_k+p^i(1-p^i)/n^i$ (sec.~A.8.3, pp.~287--288, page image of p.~288), in which the index estimate's variance is replaced by that of a single pooled proportion. Derived (corrected in the slice-34 follow-up): for the fixed-weight estimator $\hat p=\sum_jw_j\hat p_j$ and independent stock estimates with variances $\sigma_j^2$,
\[ \operatorname{Var}(\hat p_k-\hat p)=(1-w_k)^2\sigma_k^2+\sum_{j\ne k}w_j^2\sigma_j^2, \]
and with dependent estimates $\operatorname{Var}(\hat p_k-\hat p)=a^\top\Sigma a$ for $a=e_k-w$ and $\Sigma$ the covariance matrix of $(\hat p_j)$. A pooled binomial variance has no basis for an arbitrarily weighted index: with $w=(0.9,0.1)$, $p=(0.5,0.5)$ and $n=(100,100)$ the correct variance is $0.00005$ (simulated $0.000050$), and the printed formula with 200 pooled observations gives $-0.00075$, a negative variance. The per-stock variance $\sigma_j^2=p_j(1-p_j)/n_j$ holds for a \emph{count} share over $n_j$ trades only if each trade goes to the venue independently with the same probability; a \emph{value} share $\sum_\ell v_\ell X_\ell/\sum_\ell v_\ell$ with fixed trade values and independent venue indicators has variance $p_j(1-p_j)/n_{\mathrm{eff}}$, $n_{\mathrm{eff}}=(\sum v_\ell)^2/\sum v_\ell^2$ (checked: 92 effective trades out of 2{,}000 lognormal values, standard deviations 0.0478 both ways). The textbook's $n^i_k$ and $n^i$ are left undefined; its formula is an unsupported simplification, kept here only as a record of what it prints. Interpretation, checked by simulation: a value share is not a count of identical trials, and a venue choice repeated over the child orders of one parent breaks independence; with venue choices repeated in runs of 20 trades the standard deviation of a 2{,}000-trade share is 0.047 against the binomial 0.010. The comparison of a stock with an index that contains it is mechanically correlated; a leave-one-out index is the legacy notes' remedy (interpretation). \emph{Time of day.} Shares differ between the first and last hours (p.~35) and with auctions, which in European markets were mostly a primary-market monopoly; in August 2012 auctions carried 11 to 26 percent of primary-market volume across eleven European indices, and closing prices fed fund NAVs and derivatives' exchange delivery settlement prices (chap.~2, sec.~2.1, Table~2.1, pp.~117--119). Hasbrouck gives the reason auction prices carry weight: closing prices value margin positions, mutual fund shares, cash-settled derivatives and occasionally merger terms, which invites attempts to ``mark the close'' \citep[chap.~2, sec.~2.4, p.~18]{s6-hasbrouck}; the legacy ``reference price monopoly'' principle restates the primary markets' hold on auctions described here (slice~35). \emph{Proportions in general} (slice~36, from the legacy data-engineering file; a routing exception). Lehalle and Laruelle present the same estimator for any microstructure proportion: ``market shares, proportions of HFT in the market, probability to have a transaction at the bid or ask side, etc.'' (sec.~A.8, p.~284). The normal interval ``does not always work well, especially for situations with a small sample size and a proportion very close to zero or one'', and ``fails totally when the sample proportion is exactly zero or exactly one'' (pp.~285--286). Computed exactly from the binomial distribution, the nominal 95 percent Wald interval covers $p$ with probability 0.957 at $n=30$, $p=0.5$, but 0.809 at $p=0.1$ ($np=3$) and 0.454 at $p=0.02$ ($np=0.6$), and 0.908 at $n=200$, $p=0.02$ ($np=4$); Wilson's score interval covers 0.957, 0.974, 0.978 and 0.933 in the same cases; with no successes in 30 the Wald interval is the single point 0 and Wilson's is $[0, 0.114]$ (\nolinkurl{slice36-checks.py}). Lehalle and Laruelle state the rule of thumb with the true $p$; the legacy caveat applied it to $\hat p$ and recommended an exact or adjusted interval below it, which is consistent with these coverages. The legacy \emph{sampling units} (trade count, traded value, message count, time bucket, security bucket) each define a different estimand, as the count and value shares above show; ``changing the sampling unit changes the estimand'' is the legacy notes' accurate summary, and the independence needed for the binomial variance must be argued afresh for each unit.

\qffield{Practice} \emph{Volume curves} follow reference times: a linked market's opening (U.S.\ opening for European stocks; Brazil and South Africa), scheduled macroeconomic releases (the U.S.\ employment report at 14:30 CET), derivative expiries (the Eurostoxx~50 delivery settlement price, the mean index level from 11:50 to 12:00 CET on the third Friday) and reference-price rules (India's last-half-hour VWAP; Hong Kong's median of five snapshots in the last minute, with 17 percent of volume in the last 15 minutes); the patterns move with daylight saving time and market holidays (sec.~2.1.2, pp.~129--139; Total~SA, 1 January 2006 to July 2011). Rules that define a benchmark create volume around it. An average daily profile mixes such days and should not be used unfiltered for a VWAP schedule (p.~139); see \hyperref[meth:intraday-deseasonalisation]{intraday deseasonalisation}.

\qffield{Pitfalls} The legacy market-share formula matched eq.~(1.1); its ``trade-count share replaces traded value by $\mathbf 1_n$'' omitted that the denominator becomes the number of trades. Its claim that fixing auctions make a primary venue ``look more competitive'' is right only for the all-included share. Market share measured from trades is a rough prior for routing, not a measure of available liquidity (p.~38; \hyperref[def:smart-order-routing]{smart order routing}).

\qffield{See also} \hyperref[meth:fragmentation-concentration-metrics]{fragmentation and concentration metrics}; \hyperref[mod:sor-market-share-dynamics]{SOR and market-share dynamics}; \hyperref[ev:european-fragmentation-evidence]{European fragmentation evidence}; \hyperref[meth:intraday-deseasonalisation]{intraday deseasonalisation} (volume models); \hyperref[def:execution-benchmark]{execution benchmarks}.

\qffield{Further reading} Agresti and Caffo (2000) on small-sample intervals, cited by Lehalle and Laruelle, is not read. The legacy file is kept verbatim in \nolinkurl{cross-venue-routing-fragmentation.tex} in the folder \nolinkurl{notes/artefacts/market-structure-execution-quality/}. The proportion blocks merged in slice~36 are kept in \nolinkurl{notes/artefacts/lob-deep-learning/}.

\qffield{Index} market share; venue share; value share; trade-count share; ex-fixing share; fixing auctions; turnover weighting; index weighting; proportion estimator; confidence interval; binomial approximation; volume curve; reference price; exchange delivery settlement price; witching day; linked market opening; macroeconomic releases; microstructure proportions; Wald interval; Wilson interval; sampling unit

\qffield{Status} Verification: partly checked 2026-10-16 (against the text layer of a local copy of \texttt{alpha-laruelle}, 2018 World Scientific printing, chap.~1, pp.~33--39, chap.~2, pp.~117--139, and app.~A, sec.~A.8, pp.~284--288, page images of pp.~35 and~288; the ex-fixing formula derived and the estimator's dependence computed; the ex-fixing domain and the fixed-weight variance corrected in the slice-34 follow-up, 2026-10-16; slice~35: \texttt{s6-hasbrouck}, 2007 Oxford printing, p.~18; slice~36: pp.~284--286 for proportions in general, and interval coverage computed). Stability: draft. Currency: methods timeless; data and auction shares as of 2006--2013.
\end{qfmethod}

\begin{qfmethod}[Fragmentation and concentration metrics: the Fragmentation Efficiency Index, coverage, the Gini coefficient and two-sample Kolmogorov--Smirnov comparisons]\label{meth:fragmentation-concentration-metrics}
\qffield{Statement} Fragmentation exists at several levels: market operators, trading venues, order books within a venue, orders split across books, and even connection protocols; market share is the most used metric but each level needs its own \citep[chap.~1, sec.~1.1.1, pp.~33--34]{alpha-laruelle}. For shares $q_1,\dots,q_N$ summing to one, the \emph{Fragmentation Efficiency Index} is normalised entropy,
\[ \mathfrak F(q_1,\dots,q_N)=\frac{-\sum_nq_n\log q_n}{\log N}\in[0,1], \]
0 when all liquidity is in one pool and 1 when it is spread evenly (app.~A, sec.~A.1, Definition~1, pp.~247--249). \emph{Coverage} is the percentage of an index's stocks needed to make 80 percent of its turnover on a venue, 80 percent for a perfectly balanced index; the \emph{Gini coefficient} is $G=1-2B$ with $B$ the area under the Lorenz curve, equivalently half the relative mean difference, 0 for equality and near 1 for concentration (chap.~2, sec.~2.4.1, p.~159; app.~A, sec.~A.9.1, pp.~288--289). The two-sample Kolmogorov--Smirnov statistic is $D_{n,m}=\sup_x|F^X_n(x)-F^Y_m(x)|$, and equality of the two distributions is rejected at level $\alpha$ if $\sqrt{nm/(n+m)}\,D_{n,m}>k_{1-\alpha}$, the Kolmogorov quantile (sec.~A.9.2, pp.~289--291, page image of p.~291).

\qffield{Conditions} The FEI depends on the venues counted: an empty new venue lowers it from $\mathfrak F(t)$ to $\log N/\log(N+1)\cdot\mathfrak F(t)$ (about 80 percent for $N=3$, 90 for $N=6$), and Lehalle and Laruelle suggest counting a venue only above 1 percent share (p.~250). The same formula measures different things depending on the shares fed to it (trades, value, displayed depth, sessions, universes), and none measures access, connectivity or protocol costs or hidden liquidity (interpretation, consistent with p.~34). The KS test assumes independent observations from continuous distributions; venue shares across the stocks of one index are neither identically distributed nor independent of the index.

\qffield{Practice} Checked (\nolinkurl{slice34-checks.py}): shares $(0.6,0.25,0.15)$ give $\mathfrak F=0.854$, and $0.676$ with an empty fourth venue, the predicted ratio; the Lorenz and mean-difference Gini agree (0.565 on a lognormal sample of 101 stocks); Lehalle and Laruelle's Table~A.3 values follow from $G=1-2B$ (for example $B=0.18$, $G=0.64$). Their application compares SBF~120 market-share distributions on Euronext Paris and Chi-X for 2008--2012 (Tables~A.3--A.4, pp.~291--292): the KS statistics are 0.664, 0.235, 0.121 and 0.083, with $p$-values $8\times10^{-23}$, 0.003, 0.35 and 0.78. For 2010 and 2012 the test fails to reject equality at 5 percent; that does not show the distributions are the same. Checked by simulation: samples of 30 from $N(0,1)$ and $N(0.3,1.2^2)$ are not rejected in 89 percent of trials. For coverage, they report 45 percent for the FTSE~100 on the LSE, 52 percent for the AEX, 55 for the CAC~40 and 56 for the DAX, January to June 2012 (pp.~160--161).

\qffield{Pitfalls} Lehalle and Laruelle write that for 2010 and 2012 ``H0 is accepted, so the market share for the SBF 120 on Euronext Paris and Chi-X have the same distribution'' (p.~292); a non-significant test is a failure to reject, not acceptance, and with an index's worth of stocks its power against moderate differences can be low. They also report Gini values of 0.96 to 0.99 for primary-market turnover across the FTSE~100, AEX, CAC~40 and DAX (p.~160) alongside coverage near 50 percent; with the Lorenz definition of their appendix, coverage of 45 percent is hard to reconcile with a Gini of 0.99, and the chapter's construction (sorted in decreasing order, ``two times the area between the repartition curve and the bisector'', p.~160) may differ; the values are reported, not relied on. The legacy notes' Gini for ``participant activity'' and ``time segments'' extends the metric beyond the passages, which apply it to stocks and, in the appendix, to types of actors in general (p.~288).

\qffield{See also} \hyperref[meth:venue-market-share]{venue market share}; \hyperref[ev:european-fragmentation-evidence]{European fragmentation evidence} (coverage history); \hyperref[mod:sor-market-share-dynamics]{SOR and market-share dynamics}.

\qffield{Further reading} Shannon (1948), Gini (1912) and Kolmogorov (1933), cited by Lehalle and Laruelle, are not read. The legacy file is kept verbatim in \nolinkurl{cross-venue-routing-fragmentation.tex} in the folder \nolinkurl{notes/artefacts/market-structure-execution-quality/}.

\qffield{Index} fragmentation; Fragmentation Efficiency Index; FEI; entropy; venue count; coverage liquidity metric; Gini coefficient; Lorenz curve; relative mean difference; liquidity concentration; Kolmogorov--Smirnov test; two-sample test; failure to reject; test power

\qffield{Status} Verification: partly checked 2026-10-16 (against the text layer of a local copy of \texttt{alpha-laruelle}, 2018 World Scientific printing, chap.~1, pp.~33--34, chap.~2, pp.~158--161, and app.~A, secs.~A.1 and~A.9, pp.~247--250 and 288--292, page image of p.~291; FEI, Gini and KS power computed). Stability: draft. Currency: metrics timeless; data 2008--2012.
\end{qfmethod}

\begin{qfmodel}[Smart order routers and market-share dynamics: Lehalle and Laruelle's toy model of liquidity capture, learning and share feedback, and why variable liquidity loses flow]\label{mod:sor-market-share-dynamics}
\qffield{Statement} An SOR splits an order of size $V(\omega)$ across $K$ venues by a key $r=(r_1,\dots,r_K)$, $\sum_kr_k=1$, to maximise on average the liquidity it captures,
\[ C(r,\omega)=\sum_{k=1}^K\min\{r_kV(\omega),D_k(\omega)\}, \]
with $D_k(\omega)$ the liquidity venue $k$ provides when asked \citep[app.~A, sec.~A.3.1, eq.~(A.5), p.~255, page image]{alpha-laruelle}. It learns $r$ by trial and error, reducing the share of a venue that gives less liquidity than average and raising that of one that gives more (eq.~(A.7), p.~256). If a fraction $a$ of participants route with SORs, venue $k$'s share evolves as $p_k(t+1)=(1-a)p_k(t)+a\,r_k(t)$ (eq.~(A.8), p.~257). Coupled, with independent lognormal draws of $(V,D_1,\dots,D_K)$, the constraint $E[D_k/\sum_\ell D_\ell]=p_k$ (eq.~(A.10)), standard deviations proportional to the quantities through an unreported parameter $\nu$, and $E[V]\simeq\sum_kE[D_k]$, their simulation takes a 95/3/2 split, with SORs starting from equal keys, to about 50/25/25 (Figure~A.3, pp.~257--258). Shares end near their start if liquidity is steady and the routers' initial keys are close to the start, and nearer equal splits when liquidity is noisy and keys start near uniform (sec.~A.3.5, pp.~259--260).

\qffield{Notation} $\delta_{r_kV<D_k}=1$ when venue $k$ had more liquidity than was sent to it. $a$ is the share of SOR users and $\gamma$ the learning rate.

\qffield{Conditions} The book prints three slips, all in the page images: (A.5) says ``Minimize'' the expected capture, which the text maximises; (A.8) is printed $p_k(1+1)$ for $p_k(t+1)$; and (A.7) is printed with $\tfrac1K\sum_\ell\delta_\ell-\delta_k$, which would lower the key of a venue that had spare liquidity, against the text's description, so the sign is reversed (interpretation). The optimality condition (A.6) equates the venues' average indicators $E[\delta_{r_kV<D_k}]$. Derived (restated in the slice-34 follow-up): the marginal value of $r_k$ in $E[\sum_k\min\{r_kV,D_k\}]$ is $h_k(r_k)=E[V\mathbf 1\{r_kV<D_k\}]$, decreasing in $r_k$; at an optimum $h_k(r_k)=\lambda$ for every venue with $r_k>0$ (interior) and $h_k(0)\le\lambda$ for any venue left at zero (boundary), which reduces to (A.6) only when $V$ is not random. Counterexample (checked): with two equally likely scenarios, $V=1$, $D=(0.4,0.6)$ and $V=10$, $D=(6,4)$, the key $(0.5,0.5)$ equalises the unweighted indicators (both 0.5) but captures 4.95, while $(0.6,0.4)$ captures 5.4, the grid optimum. A learner consistent with this objective uses the size-weighted indicators as stochastic gradients. \emph{Illustrative variant, not a reproduction} (\nolinkurl{slice34-checks.py}): with $D_k=\mu_kX_k$, $X_k$ i.i.d.\ lognormal ($\nu=0.8$, chosen here), $V$ independent of $D$ with $E[V]=\sum\mu_k$, the means calibrated each period so that $E[D_k/\sum D]=p_k$, $a=0.3$, and projected stochastic gradient ascent with decreasing steps, the averaged keys agree with an exact solution of the conditions above (largest gap 0.001), and the optimal key is $\mu_k/\sum\mu$ by scale invariance (keys $(0.6,0.25,0.15)$ for means $(0.6,0.25,0.15)$). Because the calibration gives the largest venue $\mu_1/\sum\mu>p_1$ (0.971 for $p_1=0.95$), shares drift towards it, reaching $(0.998,0.001,0.001)$ after 60 periods: under these assumptions the coupled model concentrates rather than fragmenting. Lehalle and Laruelle's 50/25/25 path is therefore not reproduced; it depends on their unreported $\nu$, their printed update and possibly other choices, and the slice-34 claim that a corrected-sign simulation reproduced it is withdrawn (that simulation set $E[D_k]=p_k$, made $V$ depend on total liquidity, and used unweighted indicators). The model is a toy: liquidity is exogenous apart from its mean, there are no prices, fees or latency, and convergence needs the liquidity process to be i.i.d.\ or ergodic (p.~256).

\qffield{Practice} Lehalle and Laruelle's narrative version (chap.~1, sec.~1.2.2, pp.~71--73): routers adjust splits from the historical performance of their executions; ``the fraction of an order to send to a venue with a large variance will thus be minimized by a rational SOR'', and venues under-weighted long enough lose share; with shares of 76, 17 and 7 percent in their simulation, the largest market's share fell from 87 to 56 percent with medium variance and to 40 with high variance. They add that ``all studies have shown that variance is proportional with market share'', so primary markets have more variance and investors send Chi-X more than its share would imply (p.~72); that statement is reported without a source and is not checked. The legacy taxonomy of drivers (initial share, liquidity level and variance, SOR adoption $a$, reassessment speed $\gamma$) restates the model's parameters.

\qffield{Pitfalls} The legacy block, filed as an empirical finding, said that volatile fills at an alternative venue lead routers to cut its allocation and push flow back to the primary market. That direction follows from the model for a noisier alternative venue; but it is a simulation and model result, not an empirical finding, and Lehalle and Laruelle's own (unsourced) claim is that the primary is usually the noisier venue, so in their account the same mechanism erodes the primary's share.

\qffield{See also} \hyperref[def:smart-order-routing]{smart order routing}; \hyperref[meth:venue-market-share]{venue market share}; \hyperref[ev:european-fragmentation-evidence]{European fragmentation evidence}.

\qffield{Further reading} Foucault and Menkveld (2008) and Laruelle et al.\ (2011), cited by Lehalle and Laruelle, are not read. The legacy blocks are kept verbatim in \nolinkurl{cross-venue-routing-fragmentation.tex} in the folder \nolinkurl{notes/artefacts/market-structure-execution-quality/}.

\qffield{Index} smart order router; SOR; liquidity capture; routing key; learning by trading; stochastic approximation; market-share dynamics; SOR adoption; liquidity variance; reinforcement; fragmentation dynamics

\qffield{Status} Verification: partly checked 2026-10-16 (against the text layer of a local copy of \texttt{alpha-laruelle}, 2018 World Scientific printing, chap.~1, sec.~1.2.2, pp.~71--73, and app.~A, sec.~A.3, pp.~253--260, page images of pp.~255 and~257; the first-order condition derived and the dynamics simulated). Stability: draft. Currency: a stylised model of European fragmentation after 2007. Objective, update and simulation restated in the slice-34 follow-up (2026-10-16).
\end{qfmodel}

\begin{qfevidence}[European equity fragmentation after MiFID: what moved market share, which stocks fragmented, the Spanish comparison, outages and shared reference prices]\label{ev:european-fragmentation-evidence}
\qffield{Statement} Lehalle and Laruelle report, from European equity data around 2008--2013 \citep[chaps.~1--2]{alpha-laruelle}:
\begin{itemize}
\item \emph{liquidity and fragmentation}: ``the more liquid, the more fragmented'' (chap.~1, p.~37); MTFs' shares grew faster in liquid stocks, and coverage rose for all venues and indices from 2009 to 2012 but stayed lower for less liquid stocks, which they attribute to costlier hedging and short selling for high-frequency market makers, who extend to less liquid stocks once their hedges and technology are in place, down to a liquidity threshold (chap.~2, sec.~2.4.2, pp.~161--162);
\item \emph{spreads and share}: a venue's market share moves with its spread relative to the primary market's, because routers must send flow to the best price; Chi-X's spreads became tighter than the primary's and its share rose; Turquoise's share fell sharply in 2009 when the agreements under which eight market makers had provided constant liquidity ended (sec.~2.3.2, pp.~153--155);
\item \emph{resilience}: venues where trading continues as volatility rises take more of the price formation; Chi-X approached the primary market's pattern between 2008 and 2012 (sec.~2.3.3, pp.~156--158);
\item \emph{time of day}: on CAC~40 stocks (January--February 2009 and 2012) Euronext's intraday share is highest at the start and end of the day, when its book is the only one filled after the opening auction and when routers are switched on and off (sec.~2.1.3, pp.~139--140);
\item \emph{outages}: on 13 October 2010 NYSE Euronext's cash markets halted from 16:42 CET, resuming at 17:20 without market data; MTFs' midpoint books stopped and orders pegged to the primary were cancelled or paused, most large caps kept trading on MTFs near Euronext's last price, but some stocks deviated by about 2 percent, ``as if two price formation processes took place in parallel'' (sec.~2.5.2, pp.~183--185);
\item \emph{Spain}: MTFs had almost no share of IBEX~35 trading from November 2007 to January 2010 (Chi-X first exceeded 0.5 percent in August 2009), because only Iberclear, owned by the primary market BME, could clear and settle until the Titulo~V reform of June 2011 (sec.~2.5.1, p.~170).
\end{itemize}

\qffield{Conditions} The Spanish comparison is a historical contrast, not a control: comparing 15-minute Garman--Klass volatility of FTSE~100 and IBEX~35 stocks from June 2007, Lehalle and Laruelle find no significant effect of UK fragmentation on intraday volatility, and the comparison is confounded from 2010 by the sovereign-debt crisis, to which Spain was more exposed (pp.~171--174). The spread--share and resilience findings are co-movements in their figures, with Schneider Electric's five-minute data for January to July 2012 as one example (p.~156), not estimated causal effects. High-frequency traders' shares of trades (``more than 80\% \dots\ in the US, 40\% in Europe'') are reported without a source (p.~163).

\qffield{Practice} \emph{Shared references.} Lehalle and Laruelle argue that European outages showed how hard fragmented trading is without information shared among agents, and that, without a consolidated tape or trade-through rule, MTFs priced off the primary's quotes, so the primary appeared to lead (introduction, p.~5; pp.~184--185). Harris's account of consolidation is the general mechanism: in fragmented markets traders adjust to information from other segments, route to the segments with the best prices, and arbitrageurs move liquidity between segments, which consolidates the market provided information about all segments is public at low cost \citep[chap.~26, sec.~26.4, pp.~533--534]{alpha-harris}. \emph{Outages and parallel prices} (slice~36). The legacy caveat that when alternative venues keep trading while the primary's data or matching is unavailable, prices may split into parallel formation processes, is Lehalle and Laruelle's account of the Euronext cash-market halt of 13 October 2010: trading stopped at 16:42 CET and resumed at 17:20 without market data, so that for 40 minutes Euronext resembled a dark pool; most large caps kept trading on MTFs, midpoint trading stopped and orders pegged to the primary's best quotes were cancelled or paused, prices mostly stayed near Euronext's last trade, but some stocks, including EuroStoxx~50 and CAC~40 components, deviated by about 2 percent, and Euronext's later data showed its prices differing from the MTFs' \citep[chap.~2, sec.~2.5.2, pp.~183--185]{alpha-laruelle}. One episode, described from their own data; how prices on the two sides were reconciled was not quantified. Figures~2.31 and~2.32 (Vivendi and Publicis) are captioned ``October 10, 2010'', a Sunday; the text's 13 October 2010, a Wednesday, is taken as the date (pp.~185--186).

\qffield{Pitfalls} The legacy case study called Spain a ``fragmentation control case''; Lehalle and Laruelle use it as a comparison and themselves find no clear effect. Its claim that primary-market dominance persisted ``when alternative venues could not clear and settle efficiently'' is the source's explanation for Spain only. The legacy subsidised-share caveat generalised from Turquoise's 2009 episode. Venue names, shares and the regulatory framework are as of 2008--2013.

\qffield{See also} \hyperref[meth:venue-market-share]{venue market share}; \hyperref[meth:fragmentation-concentration-metrics]{fragmentation metrics}; \hyperref[mod:sor-market-share-dynamics]{SOR and market-share dynamics}; \hyperref[res:fee-adjusted-venue-choice]{fee-adjusted venue choice}; \hyperref[def:smart-order-routing]{smart order routing}; \hyperref[meth:garman-klass-volatility]{Garman--Klass volatility}; \hyperref[def:hft-strategies]{HFT strategies}.

\qffield{Further reading} The pathway's readings on fragmentation (O'Hara and Ye) and on current EU and U.S.\ rules are not read. The legacy file is kept verbatim in \nolinkurl{cross-venue-routing-fragmentation.tex} in the folder \nolinkurl{notes/artefacts/market-structure-execution-quality/}. The outage block merged in slice~36 is kept in \nolinkurl{notes/artefacts/market-structure-execution-quality/}.

\qffield{Index} market fragmentation; MiFID; MTF; Chi-X; Turquoise; BATS; coverage; high-frequency market makers; hedging cost; spread and market share; resilience to volatility; intraday market share; outage; Euronext outage 2010; Spain; IBEX 35; Iberclear; clearing and settlement barriers; consolidated tape; market consolidation; order-flow externality; Euronext outage of 2010; pegged orders

\qffield{Status} Verification: partly checked 2026-10-16 (against the text layers of local copies of \texttt{alpha-laruelle}, 2018 World Scientific printing, introduction p.~5, chap.~1, p.~37, chap.~2, pp.~139--140, 153--163 and 170--185, and \texttt{alpha-harris}, 2002 Oxford printing, chap.~26, pp.~526--534; slice~36: pp.~183--185 for the Euronext outage). Stability: draft. Currency: European markets 2007--2013; no current market structure stated.
\end{qfevidence}
